La calidad de supervisión depende de la cobertura del léxico, abreviaturas y alcance de la negación.
La identificación de una escritura no equivale a determinar el idioma: el grupo latino reúne varias
lenguas y textos mixtos. Incluso cuando una entidad se reconoce, una construcción coordinada,
una negación distante, terminología local o un matiz de incertidumbre pueden producir errores.
Una traducción literal tampoco demuestra equivalencia de criterios clínicos. El desacuerdo de las reglas
debe leerse junto con la cobertura: una regla selectiva puede aparentar mejor acuerdo al abstenerse
en los casos difíciles. Los 58 expertos no permiten validar de forma sólida todas las combinaciones
de idioma, centro y hallazgo; los CSV no identifican de manera fiable los 16 centros ni su geografía.

La selección por rareza busca cobertura de positivos y modifica la distribución respecto al universo
de estudios. Las métricas no son estimaciones de prevalencia poblacional, y el filtrado adicional por
confianza puede concentrar ciertos estilos de reporte. Las imágenes asociadas a textos con menor
cobertura tienen menos supervisión, de modo que una desigualdad lingüística puede convertirse en
desigualdad de aprendizaje visual. La ponderación preserva evidencia de baja confianza, pero también
introduce falsos negativos al tratar la ausencia de mención como un negativo débil. El filtro elimina
esa evidencia y puede dejar hallazgos con muy pocos negativos explícitos.

El riesgo de despliegue consiste en aprender correlaciones de adquisición y de anotación. Una
normalización por volumen no elimina diferencias de equipos, secuencias o poblaciones. Son necesarias
etiquetas expertas adicionales estratificadas por idioma y centro, auditoría de negación por hablantes
competentes y evaluación externa con revisión profesional. El presente ejercicio no acredita un sistema
apto para uso clínico. Los mapas de activación ofrecen indicios espaciales y no sustituyen esa revisión.
